% ============================================================
\section{Introduction}
% ============================================================

In October 2018, Lion Air Flight 610 crashed thirteen minutes after takeoff, killing 189 people. Five months later, Ethiopian Airlines Flight 302 followed the same trajectory, killing 157 more. Investigations revealed that Boeing's Maneuvering Characteristics Augmentation System (MCAS) relied on a single angle-of-attack sensor; when that sensor produced erroneous readings, MCAS repeatedly overrode pilot control inputs to push the nose down. The failure was not in the sensor, the software, or the pilots individually---it was in the interaction between layers of a constraint architecture that no single actor fully understood or controlled~\cite{jatr2019737max}.

AI system safety faces an analogous diagnostic problem in a more acute form. Multiple control nodes---users, interfaces, alignment algorithms, corporate roadmaps---shape system behavior through hierarchical control loops at different architectural layers. When these control signals conflict or feedback loops fracture, the resulting emergent hazards are not attributable to any individual component, and existing alignment paradigms lack the cybernetic vocabulary to identify where the control loop broke, which actuator was saturated, and where architectural redundancy should concentrate.

But the analogy breaks at a critical point. Systems-theoretic accident models like Leveson's STPA~\cite{leveson2012engineering, leveson2018stpa} locate causation in control-structure failures across hierarchical layers, presupposing that safe states are physically defined facts---correct airspeed, proper control-surface deflection, intact structural load paths. AI deployment lacks this property at the operational level. Terminal outcomes---death, injury, military targeting---may be uncontested, but the operational safe states that precede them (what counts as crisis-detection threshold, what counts as legitimate military use, what counts as acceptable training-data norms) are contested, layer-dependent, and retrospectively constructed. STPA decomposes safety into well-defined sub-states at each hierarchical level; AI governance must operate where that decomposition is itself disputed. The same output can be harmful or benign depending on context, audience, and the governance framework applied to it.

Technical AI research has successfully translated humanistic concepts into engineering specifications---attention became transformers, preference became reward models---but no equivalent reverse translation has given systems engineering the vocabulary to diagnose emergent hazards that originate deep within the control architecture and produce catastrophic unmodeled dynamics. This paper supplies that vocabulary by extending Leveson's STPA to a control structure whose relevant feedback is treated as structurally, not merely contingently, inadequate.

We introduce the Constraint Cascade (hereafter, the Cascade), a four-layer diagnostic taxonomy developed as a systems-safety control structure. The Cascade partitions LLM deployment into Environment (L1), Interface (L2), Algorithm \& Weights (L3), and Strategic Authority (L4), with an epistemic boundary between L2 and L3 imported as a premise from companion work~\cite{anon2026dual}. We extend Leveson's STPA to a domain where operational safe states are contested even when terminal outcomes are not, and---more sharply---where the relevant feedback can be structurally bounded. The argument follows the STPA method: we fix losses, hazards, and system-level constraints and model the control structure (\S\ref{sec:cascade}); we enumerate the Unsafe Control Actions it admits and map three loss-scenario classes---Override, Inadequate Process Model, and Compression---to the control-structure positions where the loop fractures (\S\ref{sec:patterns}); and we show the resulting control deficit cannot be closed by component-level patches, because the standard STPA remedy of improving feedback is structurally unavailable, requiring architectural redundancy and forced state disclosure instead (\S\ref{sec:responsibility}). Five cases span the range from no human agency to full deployer intent, the last prospective---a licensed clinical gatekeeper whose epistemically correct refusal the architecture cannot discharge (\S\ref{sec:cases}).

\begin{figure}[t]
\centering
\resizebox{0.95\linewidth}{!}{%
\begin{tikzpicture}[
  font=\small,
  box/.style={draw, rounded corners=3pt, line width=1pt, align=center, inner sep=6pt, text width=8.8cm},
  ca/.style={-{Latex[length=2.6mm]}, line width=1pt},
  fb/.style={-{Latex[length=2.6mm]}, line width=0.9pt, densely dashed},
  elbl/.style={font=\scriptsize, align=center, fill=white, inner sep=1.6pt},
]
% layers (top to bottom)
\node[box, fill=ccL4!35] (L4) at (0,8.1) {\textbf{L4\quad Strategic Authority}\\[1pt]\scriptsize corporate strategy \,\textbar\, sovereign override --- sets risk budget};
\node[box, fill=ccL2!45] (L2) at (0,5.4) {\textbf{L2\quad Platform Controller}\\[1pt]\scriptsize system prompts \,\textbar\, guardrails \,\textbar\, moderation \,\textbar\, RAG/context};
\node[box, fill=ccL3, text=white] (L3) at (0,1.9) {\textbf{L3\quad Automated Controller}\\[1pt]\scriptsize core agent --- process model $=$ operative state $S_t$ \textit{(hidden)}};
\node[box, fill=ccL1!55] (L1) at (0,-0.8) {\textbf{L1\quad Environment}\;\scriptsize(controlled process)\\[1pt]\scriptsize active: user (operator \& loss-bearer) \,\textbar\, passive: data-subjects};
% control actions (down, left column @ -4.0) + feedback (up, right column @ +4.0)
\draw[ca] ($(L4.south)+(-4.0,0)$) -- node[elbl]{control:\\resourcing,\\deploy mandate} ($(L2.north)+(-4.0,0)$);
\draw[fb] ($(L2.north)+(4.0,0)$) -- node[elbl]{feedback:\\evals, incidents} ($(L4.south)+(4.0,0)$);
\draw[ca] ($(L2.south)+(-4.0,0)$) -- node[elbl, pos=0.18]{control action:\\sysprompt,\\guardrail, gating} ($(L3.north)+(-4.0,0)$);
\draw[fb] ($(L3.north)+(4.0,0)$) -- node[elbl, pos=0.18]{feedback:\\trace $\tilde T_t$ \textbf{only}} ($(L2.south)+(4.0,0)$);
\draw[ca] ($(L3.south)+(-4.0,0)$) -- node[elbl]{action:\\tokens, tool calls} ($(L1.north)+(-4.0,0)$);
\draw[fb] ($(L1.north)+(4.0,0)$) -- node[elbl]{feedback:\\user response} ($(L3.south)+(4.0,0)$);
% epistemic boundary across the L2<->L3 link
\coordinate (m23) at ($(L2.south)!0.5!(L3.north)$);
\draw[ccGap, line width=1.3pt, densely dotted] ($(m23)+(-5.4,0)$) -- ($(m23)+(5.4,0)$);
\node[ccGap, fill=white, draw=ccGap, rounded corners=2pt, inner sep=2.4pt, font=\scriptsize\bfseries, align=center, text width=4.4cm] at (m23) {EPISTEMIC BOUNDARY\\ $H(S_t\mid\tilde T_t)\ge\varepsilon_{\mathrm{unrec}}^{\mathrm{LB}}$\\ {\scriptsize\mdseries dual-cert min-cut}};
% contamination: L1 -> L3, train-time, unlogged, bypassing L2 (far left)
\draw[ccRed, ca] (L1.west) to[out=160,in=200] node[elbl, text=ccRed, pos=0.72]{passive data:\\train-time,\\UNLOGGED} (L3.west);
% constitutive slow loop: L4 -> L3 weights via L3/L4 seam (far right)
\draw[ccPur, ca] (L4.east) to[out=-20,in=20] node[elbl, text=ccPur]{constitutive:\\training, resourcing\\(L3/L4 seam)} (L3.east);
\end{tikzpicture}%
}
\caption{\textbf{The Constraint Cascade as an STPA control structure.} Solid arrows (left) are control actions issued downward; dashed arrows (right) are feedback upward. L2 is a \emph{platform controller} that configures and gates L3, not a passive pipe. The red dotted line is the \emph{epistemic boundary} of the companion result~\cite{anon2026dual}: L3's operative state $S_t$ is non-recoverable from the trace $\tilde T_t$, so L2's process model of L3 is structurally incorrect---an inadequate-process-model hazard \emph{between two controllers}, the case a state-observable (Mealy/Moore) abstraction excludes. The min-cut structural budget $\varepsilon_{\mathrm{unrec}}^{\mathrm{LB}}$ is the one quantity an external reviewer can still compute, used here as a feedback-adequacy meter. The contamination variant enters on the unlogged train-time edge L1$\to$L3 (red), bypassing L2 entirely; the constitutive slow loop L4$\to$L3 (purple) sets L3's control algorithm through training and resourcing at the L3/L4 seam. Override, Inadequate Process Model, and Compression are the loss-scenario classes localized to these positions (\S\ref{sec:patterns}).}
\label{fig:cascade}
\end{figure}
